\subsection{高原、海拔与性健康}

高原是常被忽略的性健康影响因素。低氧、低气压、干燥与低温共同作用，会通过血管、内分泌与心理多重途径影响性欲与性功能。对进藏旅行者、高原作业人员和高原常住居民而言，了解这些规律既能避免不必要的恐慌，也能提前做好防护。

\vspace{0.5em}
\noindent\textbf{低氧如何影响性功能}：

\begin{itemize}
	\item \textbf{血管与血流}：勃起依赖阴茎动脉的充分扩张与海绵体血流充盈。急性高原暴露下，交感神经兴奋、外周血管收缩，可能使勃起硬度下降、维持困难——通常为\textbf{暂时性}，随适应（acclimatization）而改善
	\item \textbf{内分泌反应}：多项研究显示，高原低氧环境可致\textbf{睾酮水平短期下降}、促性腺激素分泌改变；一般在适应数周后部分回升，但个体差异大
	\item \textbf{红细胞增多与血液黏稠}：长期高原居住者红细胞代偿性增多、血液黏滞度升高，理论上不利于微循环——对已有心血管基础病者，性活动时的心脏负荷更需评估
	\item \textbf{性欲与情绪}：高原反应（头痛、失眠、食欲减退、乏力）本身即明显降低性欲；\textbf{高原失眠}与夜间周期性呼吸（睡眠呼吸紊乱）进一步耗损精力
\end{itemize}

\vspace{0.5em}
\noindent\textbf{高原性生活的实用建议}：

\begin{itemize}
	\item \textbf{抵达后先适应，不急于性活动}：一般建议抵达高原后 1--3 天内避免剧烈体力活动（包括性活动），让身体先完成初步适应；有严重高原反应者应推迟至症状消退
	\item \textbf{控制强度，量力而行}：高原性活动的心肺负荷高于平原——把"能轻松完成日常活动无气促"作为参考标准，避免过度劳累；出现明显气促、心悸、胸痛应立即停止
	\item \textbf{心血管病患者需个体化评估}：冠心病、高血压、心衰患者前往高原前应咨询心内科医生，评估海拔耐受与性活动安全性（参考"心血管疾病与性"相关章节的体能自评方法）
	\item \textbf{充分补水、避免饮酒}：高原干燥加低氧，脱水与酒精都会加重缺氧与勃起困难；酒精助"兴"是错觉，实为抑制
	\item \textbf{注意保暖与润滑}：高原寒冷、空气干燥，可能加重阴道干涩与性交不适——合理使用润滑剂，注意保温
\end{itemize}

\vspace{0.5em}
\noindent\textbf{高原与生育}：低氧环境对精子质量的影响尚无一致结论——部分研究提示高原暴露期间精子浓度与活力下降，回到平原后多可恢复；备孕夫妇如有条件，可避开刚上高原的适应期安排受孕。孕妇前往高原应格外谨慎，需产科医生评估。

\begin{tcolorbox}[colback=pink!5!white,colframe=red!50!black,title=贴心小叮咛]
高原对性功能的影响\textbf{多数是暂时的}——身体的优先级是"先活下来"，性欲下降是身体的合理节能策略，不必惊慌或自责。给身体几天适应时间，规律作息、充分休息，大多数情况会自行好转；持续存在或伴胸痛、呼吸困难，请及时就医。
\end{tcolorbox}
